\subsection{Our contribution}
The work of~\cite{streamexpdec} initiated the study of expander decompositions in the streaming setting, and consequently raised the question of whether a dependence in space on the sparsity parameter~$\phi$ is inherent. In this paper, we move towards settling the streaming complexity of expander decompositions by attacking the problem on two fronts: (1) we give a nearly optimal algorithm for ``one-level'' expander decomposition that avoids the sparsity dependence, and (2) we show that computing a ``repeated'' expander decomposition, commonly used in applications, cannot avoid such dependence.
\subsubsection*{Upper bound}
We give an $\otil(n)$ space algorithm for computing a $(O(\phi \log n), \phi)$-ED in dynamic streams. Specifically, we show that a ``universal'' sketch consisting of $\otil(n)$ random linear measurements of the incidence matrix can be decoded into a $(O(\phi \log n), \phi)$-ED for any $\phi$: the sketch is independent of the sparsity~$\phi$.

\begin{mainthsimple}[ED algorithm -- exponential time decoding]
	\label{cor:main}
	Let $G=(V,E)$ be a graph given in a dynamic stream. Then, there is an algorithm that maintains a linear sketch of $G$ in $\otil(n)$ space. For any $\phi \in (0,1)$, the algorithm decodes the sketch to compute a $(O(\phi \log n), \phi)$-ED of $G$ with high probability, in $\otil(n)$ space and $2^{O(n)}$ time.
\end{mainthsimple}

\noindent
We note that at least $\Omega(n\log n)$ space is needed for any small enough $\phi$: for example, if the input graph is a matching of size $n/10$, say, its ED gives a $1-O(\phi \log n)$ fraction of the matching edges. 

The decoding time of our sketch can be made polynomial, at the expense of some loss in the quality of the expander decomposition (similarly to~\cite{streamexpdec}), but keeping the space independent of the sparsity $\phi$.

\begin{mainthsimple}[ED algorithm -- polynomial time decoding]
	\label{cor:mainpoly}
	Let $G=(V,E)$ be a graph given in a dynamic stream. Then, there is an algorithm that maintains a linear sketch of $G$ in $n^{1+o(1)}$ space. For any $\phi \in (0,1)$, the algorithm decodes the sketch to compute a $(\phi \cdot n^{o(1)}, \phi)$-ED of $G$ with high probability, in $n^{1+o(1)}$ space and $\poly(n)$ time.
\end{mainthsimple}

\noindent
In this case we are off by subpolynomial factors in both quality and space complexity as compared to the optimal ones. The actual theorem that we prove allows one to trade-off the loss in quality and increase in space.

\subsubsection*{Lower bound}
Most algorithmic applications of EDs, including the ones mentioned above, do not use just one ED of the input graph. Rather, they use an ED sequence obtained by repeatedly computing an ED of the inter-cluster edges from the previous level. This can be done in two natural ways: by contracting the clusters of an ED (we call this variant CED, for “contraction”), or by removing the intra-cluster edges without changing the vertex set (we call this variant RED, for “removal”). These approaches lead to different results and are used for different applications (e.g.,~\cite{boundlink,mincutjli} contract the clusters, and~\cite{stspectral, maxflow} recurse on inter-cluster edges). In the sequential setting, both CEDs and REDs can be obtained straightforwardly given an ED algorithm. However, this is not so obvious in the streaming model.

One the one hand, one should be able get a sparsity-independent algorithm for computing CEDs in dynamic streams via \cref{cor:main} or \cref{cor:mainpoly}: observe that contracting the vertices of a sparsifier of the input graph $G$  gives a sparsifier of the graph obtained by contracting vertices in~$G$, so the idea would be to maintain an independent copy of our algorithm for each of the $O(\log n)$ levels and contracting vertices in the sparsifier based on the decomposition of the previous level. On the other hand, we show a space lower bound for computing REDs, even in insertion only streams, showing that a dependence on $1/\phi$ is necessary.

\begin{mainthsimple}[RED lower bound]
	\label{th:lbsimple}
	Let $\epsilon,\phi \in (0,1)$ such that $1/n \ll \phi \ll 1/\pylog n$ and $ \epsilon =\otil( \phi)$. Any streaming algorithm that with constant probability computes at least two levels of an $(\epsilon,\phi)$-RED requires \smash{$\widetilde{\Omega}(n/\phi)$} bits of space.
\end{mainthsimple}

\noindent
The result seems to challenge the intuition that EDs become weaker as $\phi$ and $\epsilon$ decrease (note that when $\phi$ is, say, $1/n^2$, an ED can simply consist of connected components). The questions of~(1)~whether this bound can be improved,~(2)~how it scales with the number of levels of RED we compute, and~(3)~whether there are algorithms matching such bounds, remain open.